\documentclass[11pt]{article}

\usepackage[margin=1in]{geometry}
\usepackage{amsmath,amssymb}
\usepackage{booktabs}

\title{Denoised Waveform-Shape Metric Filtering Methods}
\author{AngioEye pipelines}
\date{}

\begin{document}
\maketitle

\section{Overview}

The denoised waveform-shape metric pipelines apply filtering to arterial
per-segment pulse waveforms before computing arterial waveform metrics.  Each
arterial pulse is represented as
\[
    y_i(t)=v[t,\mathrm{beat},\mathrm{branch},\mathrm{radius}],
\]
where \(i\) indexes one \((\mathrm{beat},\mathrm{branch},\mathrm{radius})\)
realization and \(t\) indexes samples within the beat.

All three methods first reject pulses with no finite samples, too few finite
samples, or near-zero variance.  Missing samples in valid pulses are linearly
interpolated for the filtering computation, and the original missing-sample
mask can be restored after filtering.  By default, these new methods use the
bandlimited arterial segment signal as input and use the filtered result as
the arterial metric signal.

\section{Cross-Correlation Non-Local Averaging}

Pipeline:
\[
    \texttt{waveform\_shape\_metrics\_denoised\_correl}.
\]

This method treats every valid arterial segment pulse as a noisy realization
of a shared pulse morphology.  For each target pulse, it finds other pulses
with similar shape, phase-aligns them by circular lag, optionally rescales
their amplitude to the target, averages them into a local template, and blends
that template back into the target pulse.

For a filled pulse \(\tilde y_i\), the normalized matching signal is
\[
    z_i =
    \frac{\tilde y_i-\bar y_i\mathbf{1}}
         {\|\tilde y_i-\bar y_i\mathbf{1}\|_2}.
\]
For each candidate neighbor \(j\), the algorithm searches circular lags
\[
    \ell \in \{-L_{\max},\ldots,L_{\max}\},
    \qquad
    L_{\max} =
    \min\left(\left\lfloor \rho_{\max}T \right\rceil,
    \left\lfloor T/2 \right\rfloor - 1\right),
\]
with \(\rho_{\max}=\texttt{correlation\_max\_lag\_fraction}\).  The retained
correlation and lag are
\[
    c_{ij}=\max_{\ell} z_i^\top \operatorname{roll}(z_j,\ell),
    \qquad
    \ell_{ij}=\arg\max_{\ell} z_i^\top \operatorname{roll}(z_j,\ell).
\]
Only neighbors with \(c_{ij}\geq\texttt{correlation\_min\_corr}\) are kept.

The neighbor weight is
\[
    w_{ij}
    =
    \max(c_{ij},0)^p
    \exp\left[-\frac{1}{2}
        \left(\frac{\ell_{ij}}{\sigma_{\ell}T}\right)^2\right]
    s_{ij},
\]
where \(p=\texttt{correlation\_corr\_power}\),
\(\sigma_{\ell}=\texttt{correlation\_lag\_sigma\_fraction}\), and \(s_{ij}\)
is an optional branch/radius locality factor.  The largest
\(\texttt{correlation\_top\_k\_neighbors}\) neighbors are retained.

If amplitude matching is enabled, an aligned neighbor is transformed to the
target mean and standard deviation:
\[
    \hat y_{j\rightarrow i}
    =
    \bar y_i+
    \left(\operatorname{roll}(\tilde y_j,\ell_{ij})-\bar y_j\right)
    \frac{\sigma_i}{\sigma_j}.
\]
The local template is
\[
    m_i(t)=
    \frac{\lambda_0\tilde y_i(t)+
    \sum_{j\in\mathcal{N}_i}w_{ij}\hat y_{j\rightarrow i}(t)}
    {\lambda_0+\sum_{j\in\mathcal{N}_i}w_{ij}},
\]
where \(\lambda_0=\texttt{correlation\_self\_weight}\).  The final filtered
pulse is
\[
    x_i(t)=(1-\alpha)\tilde y_i(t)+\alpha m_i(t),
    \qquad
    \alpha=\texttt{correlation\_blend\_alpha}.
\]
The result is clipped to the finite range of the original pulse and, by
default, the original \texttt{NaN} mask is restored.

\begin{tabular}{@{}ll@{}}
\toprule
Main default & Value \\
\midrule
\texttt{correlation\_top\_k\_neighbors} & 12 \\
\texttt{correlation\_min\_corr} & 0.65 \\
\texttt{correlation\_corr\_power} & 3.0 \\
\texttt{correlation\_blend\_alpha} & 0.75 \\
\texttt{correlation\_self\_weight} & 0.25 \\
\texttt{correlation\_max\_lag\_fraction} & 0.12 \\
\texttt{correlation\_lag\_sigma\_fraction} & 0.06 \\
\bottomrule
\end{tabular}

\section{Graph-Laplacian Pulse Denoising}

Pipeline:
\[
    \texttt{waveform\_shape\_metrics\_denoised\_laplace}.
\]

This method builds a graph whose vertices are valid arterial pulses.  Edges
connect pulse realizations with similar phase-aligned morphology.  The
denoised ensemble solves the graph Tikhonov problem
\[
    \min_X \|X-Y\|_F^2+\gamma\operatorname{tr}(X^\top L X),
\]
where rows of \(Y\) and \(X\) are pulse realizations, columns are time samples,
\(L\) is the pulse-graph Laplacian, and
\(\gamma=\texttt{laplacian\_gamma}\).

The method first constructs a reference pulse from the median valid pulse,
falling back to the mean if needed.  Each pulse is circularly aligned to that
reference using the same bounded lag search idea as the correlation method.
After alignment, pulse similarities are computed at zero lag.

A directed candidate edge \(i\rightarrow j\) is retained when the aligned
correlation satisfies
\[
    c_{ij}\geq\texttt{laplacian\_min\_corr}.
\]
Its weight is
\[
    w_{ij}
    =
    \max(c_{ij},0)^p
    \exp\left[-\frac{1}{2}
        \left(\frac{\ell_i-\ell_j}{\sigma_{\ell}T}\right)^2\right]
    s_{ij},
\]
where \(p=\texttt{laplacian\_corr\_power}\), \(\ell_i\) and \(\ell_j\) are
the reference-alignment lags, and \(s_{ij}\) optionally penalizes branch/radius
index distance.  Only the largest
\(\texttt{laplacian\_top\_k\_neighbors}\) directed weights are retained per
pulse.  The graph is symmetrized by
\[
    W_{ij}=\max(W^{\mathrm{directed}}_{ij},
                W^{\mathrm{directed}}_{ji}),
\]
and the Laplacian is
\[
    L=D-W,\qquad D_{ii}=\sum_j W_{ij}.
\]

The objective gives the linear system
\[
    (I+\gamma L)X=Y.
\]
The implementation adds a small diagonal
\(\texttt{laplacian\_self\_jitter}I\) for numerical stability and solves the
dense system.  After solving, the result is blended with the aligned input:
\[
    X_{\mathrm{final}}=(1-\alpha)Y+\alpha X,
    \qquad
    \alpha=\texttt{laplacian\_blend\_alpha}.
\]
Rows are then shifted back to their original phase.  By default the solve is
performed on standardized pulse shapes and restored to each pulse's original
mean and standard deviation afterward.

\begin{tabular}{@{}ll@{}}
\toprule
Main default & Value \\
\midrule
\texttt{laplacian\_gamma} & 0.75 \\
\texttt{laplacian\_top\_k\_neighbors} & 6 \\
\texttt{laplacian\_min\_corr} & 0.70 \\
\texttt{laplacian\_corr\_power} & 3.0 \\
\texttt{laplacian\_self\_jitter} & \(10^{-8}\) \\
\texttt{laplacian\_blend\_alpha} & 1.0 \\
\texttt{laplacian\_branch\_sigma} & 2.0 \\
\texttt{laplacian\_radius\_sigma} & 2.0 \\
\bottomrule
\end{tabular}

\section{Joint Pulse-Graph and Temporal-Graph Denoising}

Pipeline:
\[
    \texttt{waveform\_shape\_metrics\_denoised\_joint}.
\]

This method extends the graph-Laplacian denoiser by adding a temporal graph
over within-pulse samples.  The pulse graph encourages similar pulse
realizations to have similar denoised shapes, while the temporal graph
suppresses high-frequency variation inside each pulse.

The joint objective is
\[
    \min_X
    \|X-Y\|_F^2
    +\gamma_p\operatorname{tr}(X^\top L_p X)
    +\gamma_t\operatorname{tr}(X L_t X^\top),
\]
where \(L_p\) is the pulse-similarity graph Laplacian,
\(L_t\) is the temporal Laplacian,
\(\gamma_p=\texttt{joint\_pulse\_gamma}\), and
\(\gamma_t=\texttt{joint\_temporal\_gamma}\).

The pulse graph \(L_p\) is built from reference-aligned pulses using the same
neighbor selection and weighting strategy as the graph-Laplacian method, with
the \texttt{joint\_*} parameters.  The temporal graph \(L_t\) connects adjacent
time samples:
\[
    (L_t)_{k,k}\mathrel{+}=1,\quad
    (L_t)_{k+1,k+1}\mathrel{+}=1,\quad
    (L_t)_{k,k+1}\mathrel{-}=1,\quad
    (L_t)_{k+1,k}\mathrel{-}=1.
\]
When \texttt{joint\_temporal\_cycle} is enabled, the first and last time
samples are also connected, making the temporal graph cyclic.

The optimality condition is a Sylvester-type system:
\[
    (I+\gamma_p L_p)X+\gamma_t X L_t=Y.
\]
The implementation solves it spectrally.  With
\[
    A_p=I+\gamma_p L_p+\epsilon I,
    \qquad
    A_p=U_p\Lambda_pU_p^\top,
    \qquad
    L_t=U_t\Lambda_tU_t^\top,
\]
where \(\epsilon=\texttt{joint\_self\_jitter}\), the solution is
\[
    X =
    U_p
    \left[
        \frac{U_p^\top Y U_t}
             {\lambda_p+\gamma_t\lambda_t}
    \right]
    U_t^\top.
\]
The denominator is evaluated for every pulse-eigenvalue and temporal-eigenvalue
pair.  The reported condition number is the ratio of the largest and smallest
absolute denominator.

The solved ensemble is blended with the aligned input using
\[
    X_{\mathrm{final}}=(1-\alpha)Y+\alpha X,
    \qquad
    \alpha=\texttt{joint\_blend\_alpha}.
\]
Rows are shifted back to the original pulse phase, restored to the original
pulse scale when enabled, optionally clipped, and finally masked with the
original \texttt{NaN} locations.

\begin{tabular}{@{}ll@{}}
\toprule
Main default & Value \\
\midrule
\texttt{joint\_pulse\_gamma} & 0.75 \\
\texttt{joint\_temporal\_gamma} & 0.03 \\
\texttt{joint\_top\_k\_neighbors} & 6 \\
\texttt{joint\_min\_corr} & 0.70 \\
\texttt{joint\_corr\_power} & 3.0 \\
\texttt{joint\_self\_jitter} & \(10^{-8}\) \\
\texttt{joint\_blend\_alpha} & 1.0 \\
\texttt{joint\_temporal\_cycle} & \texttt{True} \\
\texttt{joint\_branch\_sigma} & 2.0 \\
\texttt{joint\_radius\_sigma} & 2.0 \\
\bottomrule
\end{tabular}

\section{Common Diagnostics and Fallback Codes}

All three methods record denoising diagnostics under
\[
    \texttt{artery/by\_segment/denoising}.
\]
Common diagnostics include finite-sample fraction, original-versus-filtered
correlation, status code, and mean absolute pointwise change.  The graph-based
methods also record phase-alignment lag, phase-alignment correlation, graph
neighbor statistics, graph degree, energy values, edge counts, valid pulse
count, and system condition number.

\begin{tabular}{@{}ll@{}}
\toprule
Status code & Meaning \\
\midrule
0 & Pulse was denoised. \\
1 & Pulse contained no finite samples. \\
2 & Pulse had too few finite samples. \\
3 & Pulse had near-zero variance. \\
4 & No eligible neighbors were found. \\
5 & Graph or joint solve failed. \\
\bottomrule
\end{tabular}

Status code 5 is only used by the graph-Laplacian and joint methods.  In
fallback cases, valid pulses are returned without the corresponding graph
smoothing step rather than being replaced by invalid values.

\end{document}
